% Programming with InitechBase, Chapter 15: Reports from Programs (AN-16;
% annex Sec 8.2 P12).  Panel replayed (examples/p15-pagerpt.dot); the
% program PAGERPT.PRG of the sample diskette.
\chapter{Reports from Programs}
\Gates{test-pubs-examples, test-interp-flow, test-interp-proc}
A report form lays out a report in columns.  When a report needs a layout of
its own -- lines built from several fields, a heading of its own, a count of
lines -- a program prints it with \cmd{?}.  This chapter prints the staff
roster in numbered pages.

\begin{ChapterCovers}
\item Counting lines
\item A heading on every page
\item Page numbers
\end{ChapterCovers}

\ATEHead{Preparing for This Chapter}
Start InitechBase with a fresh copy of the sample diskette, and type
\type{USE STAFF}.

\ATEHead{Pages of Lines}
A page holds a fixed number of lines.  The program counts the lines it has
printed on the page in a memory variable; when the page is full, it prints the
heading of a new page and begins the count again.  Starting the count at the
full value makes the program print the heading before the first line, too.

\ATEHead{The Roster in Pages}
\file{PAGERPT.PRG} prints six employees to a page, so that the example fits this
book; for a printed page of 55 lines, set \texttt{MLINES} to 55.
\ProgramListing{PAGERPT.PRG}
\begin{Steps}
\item Type \type{DO PAGERPT}.
\end{Steps}
\DotExample{examples/p15-pagerpt}
\SideHead{MLINE}counts the lines of the current page; \texttt{MPAGE} the pages.

\SideHead{HEADING}is the procedure that begins a page: an empty line, the title
with the page number, the column headings, and the count set back to zero.

\SideHead{The end}After the loop, a closing line tells the reader that the
roster is complete.

\ATEHead{Where Do You Go From Here?}
Chapter~16 calculates with dates.

\ATEHead{Summary}
\begin{enumerate}
\item Count the lines printed on a page in a memory variable.
\item Print the heading when the count reaches the length of the page.
\item Keep the page number in a memory variable and print it in the heading.
\end{enumerate}
